\bookchapter{Bab 8: Komunikasi: SMS, Telefon dan Sosial}

\chapterintro{Komunikasi ialah misi paling asas dan paling teras bagi telefon. Kini ia telah berkembang daripada panggilan dan SMS ringkas menjadi ekosistem kompleks yang merangkumi pemesejan segera, rangkaian sosial dan e-mel.}

\booksection{Fungsi Lanjutan Aplikasi Telefon}

Aplikasi telefon pada telefon moden sudah lama melampaui fungsi pad dail ringkas; ia mengintegrasikan alat berkuasa seperti rakaman panggilan, sekatan gangguan dan pengecaman pintar.

\booksubsection{Rakaman panggilan (perhatikan pematuhan undang-undang):}

\begin{itemize}
  \item Fungsi: Merakam kandungan panggilan penting untuk nota kerja, temu bual, janji pelanggan atau simpanan bukti bagi pertikaian.
  \item Petua undang-undang: Di China, syarat untuk menggunakan rakaman panggilan sebagai bukti ialah mendapatkan persetujuan pihak yang satu lagi, atau menggunakannya untuk melindungi hak dan kepentingan sah anda sendiri.
  \item Rakaman satu pihak vs persetujuan kedua-dua pihak: Dari segi undang-undang, rakaman yang diperoleh "tanpa persetujuan pihak yang satu lagi" tidak semestinya tidak sah. Menurut tafsiran undang-undang berkaitan Mahkamah Rakyat Tertinggi, selagi cara perolehan rakaman tidak melanggar hak dan kepentingan sah pihak lain secara teruk, tidak melanggar peruntukan larangan undang-undang dan tidak bercanggah dengan tatasusila awam, ia masih boleh dikemukakan sebagai bukti. Tetapi sesetengah mahkamah di wilayah tertentu lebih ketat menerima rakaman "senyap-senyap" dalam perbicaraan sivil; kes jenayah pula biasanya memerlukan semakan kesahan yang lebih ketat.
  \item Perbezaan antara wilayah: Mahkamah di Beijing, Shanghai, Guangdong dan tempat lain agak bertolak ansur dengan rakaman satu pihak sebagai bukti sivil; manakala rakaman yang melibatkan rahsia perniagaan atau privasi peribadi perlu sangat berhati-hati di mana-mana wilayah, dan mungkin dikecualikan kerana melanggar hak privasi.
  \item Cara beroperasi: Android biasanya mempunyai butang "Rakam" yang jelas pada skrin panggilan; sistem iOS tidak menyediakan fungsi rakaman panggilan asli, jadi perlu guna aplikasi pihak ketiga atau fon Bluetooth yang ada fungsi rakaman.
\end{itemize}

\begin{tcbwarn}{Batasan Wajib Penggunaan Rakaman}
\begin{itemize}
  \item Rakaman tidak boleh digunakan untuk ugutan, penyebaran atau penghinaan terbuka terhadap orang lain.
  \item Kandungan yang melibatkan privasi pihak ketiga (contohnya perbualan orang lain yang dirakam secara tidak sengaja) tidak boleh disebarkan sesuka hati.
  \item Untuk situasi penting, disyorkan memberitahu secara lisan pada awal panggilan bahawa "panggilan ini akan dirakam" — ini mematuhi peraturan dan menjadikan rakaman lebih berkuasa sebagai bukti.
\end{itemize}

\end{tcbwarn}

\booksubsection{Kandungan mendalam tentang penanganan panggilan gangguan:}

\begin{itemize}
  \item Prinsip: Aplikasi telefon merujuk pangkalan data nombor yang besar melalui rangkaian, lalu memaparkan "panggilan gangguan", "panggilan penipuan" atau nama syarikat apabila panggilan masuk.
  \item Cara mengaktifkan: Pada Android, masuk Aplikasi Telefon > Tetapan > Sekatan Gangguan; pada iOS, masuk Tetapan > Telefon > Sekatan dan Pengecaman Panggilan. Untuk laluan suis terperinci setiap jenama dan item yang disyorkan, lihat jadual perbandingan jenama dalam Bab 13 "Alat Sistem Praktikal".
  \item Platform penanda utama: Penandaan nombor di China disediakan terutamanya oleh tiga penyedia perkhidmatan data — Tencent Mobile Manager, 360 Mobile Guard dan Dianhuabang (Tapir). Fungsi paparan "nama pemilik" dan "sekatan gangguan" oleh operator (China Mobile/China Unicom/China Telecom) juga sering menggunakan data ini. Pangkalan penanda pengeluar telefon (Huawei, Xiaomi, OPPO, vivo) kebanyakannya bekerjasama dengan penyedia perkhidmatan di atas.
  \item Apa perlu buat jika ditanda salah: Jika nombor kerja normal anda ditanda secara salah sebagai "penghantaran makanan", "kurier" atau "panggilan gangguan", anda boleh membuat rayuan melalui saluran berikut: Dianhuabang: Lawati platform rayuan nombor (haoma.tk.cn) dan serahkan bukti identiti.
  \item 360: Serahkan melalui fungsi "Rayuan Nombor" dalam aplikasi 360 Mobile Guard.
  \item Tencent: Serahkan melalui halaman "Rayuan Nombor" dalam mini program "Tencent Mobile Manager".
  \item Tiga operator utama: Pengguna China Mobile boleh hantar SMS "KTFSR" ke 10086 untuk mengaktifkan fungsi anti-gangguan; pengguna China Telecom boleh ikuti akaun rasmi "Tianyi Anti-Harassment"; pengguna China Unicom boleh tetapkan dalam "Unicom Mobile Manager".
\end{itemize}

Anti-gangguan lanjutan: Aktifkan "Sekatan Panggilan Frekuensi Tinggi" operator (contohnya Had Panggilan Frekuensi Tinggi China Mobile) untuk menyekat secara automatik panggilan keluar suara AI.

Keupayaan biasa sekatan bawaan telefon: Pengecaman Panggilan Masuk / Sekatan Pintar: Menanda dan menyekat secara automatik gangguan, penipuan dan promosi iklan berdasarkan pangkalan data nombor.

Senarai hitam / senarai putih: Sekat secara manual nombor gangguan yang diketahui; masukkan kenalan serta nombor kurier dan penghantaran makanan ke dalam senarai putih untuk mengelak salah sekat.

Penapis kata kunci SMS: SMS yang mengandungi kata kunci seperti "menang hadiah", "pinjaman", "klik pautan" akan dimasukkan secara automatik ke dalam kotak sekatan.

iOS "Senyapkan Pemanggil Tidak Dikenali": Panggilan daripada nombor yang tiada dalam kenalan hanya akan disenyapkan dan masuk ke peti suara tanpa mengganggu anda; sesuai untuk pengguna yang kerap menerima panggilan gangguan, tetapi berhati-hati agar tidak terlepas panggilan penting daripada nombor tidak dikenali.

Kerjasama mod Jangan Ganggu: Pada waktu malam atau semasa mesyuarat, benarkan hanya panggilan daripada kenalan, digabungkan dengan suis berjadual — kualiti tidur terus bertambah baik.

Cara perisian praktikal bekerjasama: Utamakan sekatan bawaan sistem, kemudian tambah National Anti-Fraud Center dan perkhidmatan anti-gangguan operator; untuk pembersihan boleh pilih alat ringan seperti Qingzhuo, tetapi bagi kategori yang sama (keselamatan/pembersihan) satu set sahaja sudah cukup, lihat Bab 13 untuk butiran.

\booksubsection{Apa Itu VoLTE/VoNR dan Cara Mengaktifkan:}

\begin{itemize}
  \item VoLTE (Voice over LTE): Teknologi panggilan suara berasaskan rangkaian 4G; semasa panggilan ia tidak lagi turun ke 2G/3G, anda boleh melayari internet sambil menelefon, kualiti suara lebih jelas daripada panggilan tradisional dan sambungan lebih pantas.
  \item VoNR (Voice over NR): Panggilan suara berasaskan 5G, iaitu evolusi generasi akan datang selepas VoLTE. Kelewatan lebih rendah, kualiti suara lebih baik, dan menyokong panggilan video HD. Memerlukan ketiga-tiga: telefon 5G + pelan 5G + liputan 5G di kawasan anda.
  \item Cara mengaktifkan: Pada iOS, aktifkan "VoLTE" dalam "Tetapan > Rangkaian Selular > Pilihan Data Selular > Suara dan Data"; pada Android, cari suis "Panggilan HD VoLTE" dalam "Tetapan > Rangkaian Mudah Alih > Data Mudah Alih > Lanjutan". Operator biasanya mengaktifkannya secara lalai; jika anda mendapati rangkaian terputus semasa panggilan dan tidak boleh melayari internet, kemungkinan besar VoLTE tidak diaktifkan.
\end{itemize}

\begin{tcbinfo}{Tahukah Anda}
Jika kelajuan internet telefon anda tiba-tiba perlahan atau terputus semasa panggilan, kemungkinan besar VoLTE/VoNR tidak diaktifkan dan ia turun ke rangkaian 2G/3G. Selepas diaktifkan, anda boleh “tonton video sambil menelefon”.

\end{tcbinfo}

\booksection{SMS dan iMessage/RCS}

SMS tradisional tidak luput; ia masih saluran yang boleh dipercayai untuk menerima kod pengesahan, notis bank dan maklumat rasmi.

\begin{itemize}
  \item iMessage: Mesej yang dihantar melalui internet antara peranti Apple, dipaparkan sebagai gelembung biru. Menyokong gambar dan video berkualiti tinggi, resit dibaca dan sebagainya, sepenuhnya percuma dan disulitkan hujung-ke-hujung.
  \item RCS: Protokol komunikasi generasi akan datang yang didorong oleh operator global dan Google. Fungsinya serupa dengan iMessage atau WeChat, tetapi berasaskan nombor telefon.
  \item Interoperabiliti iMessage dan RCS: Bermula dengan iOS 18, Apple mula menyokong standard RCS, bermakna iPhone dan telefon Android boleh menghantar gambar HD, resit dibaca dan petunjuk sedang menaip melalui RCS, menamatkan jurang pengalaman "gelembung hijau". Tetapi syaratnya kedua-dua operator mesti mengaktifkan RCS dan versi sistem telefon cukup baharu.
\end{itemize}

\booksubsection{Keadaan pelaksanaan RCS di China:}

\begin{itemize}
  \item China Mobile: Melaksanakan RCS lebih awal; sesetengah wilayah sudah menyokong perkhidmatan "Mesej 5G", boleh menghantar mesej kaya dan kad interaktif. Memerlukan sokongan peranti telefon dan pengaktifan di kedai perkhidmatan operator.
  \item China Telecom: Menyediakan perkhidmatan melalui "Telecom 5G Message"; sesetengah pelan menyokongnya secara lalai, manakala pengguna lama perlu menghantar arahan SMS tertentu untuk mengaktifkannya.
  \item China Unicom: Melaksanakan Mesej 5G di bandar-bandar utama, tetapi liputan peranti tidak sekata; banyak model Android masih perlu mengaktifkan suis "Maklumat Dipertingkat" atau "Mesej 5G" secara manual dalam aplikasi SMS.
  \item Pengalaman sebenar: Kelajuan promosi RCS di China lebih perlahan daripada jangkaan, terutamanya dipengaruhi oleh dominasi ekosistem WeChat dan penyesuaian peranti yang tidak seragam. Komunikasi harian masih bergantung pada WeChat; RCS lebih banyak memainkan peranan "pemancar rasmi" (seperti notis bank dan mesej kerajaan).
\end{itemize}

\begin{tcbwarn}{Batasan Keselamatan Kod Pengesahan}
\begin{itemize}
  \item Jangan teruskan kod pengesahan: Sesiapa sahaja yang meminta anda "sebut kod pengesahan" atau "majukan SMS" mungkin seorang penipu, termasuk orang yang mendakwa diri sebagai "khidmat pelanggan", "kurier" atau "polis, pendakwa dan mahkamah". Kod pengesahan = beg wang anda.
  \item Elak SMS pemancingan data: Berhati-hati dengan pautan dalam SMS seperti "ETC tamat tempoh", "kad perubatan dibekukan", "kurier hilang tuntutan ganti rugi"; institusi rasmi tidak akan meminta nombor kad bank atau kata laluan melalui pautan SMS.
  \item Pengecaman nombor: SMS rasmi bank biasanya datang daripada nombor pendek yang bermula dengan 95 atau 1069 dengan tandatangan korporat; "notis bank" yang dihantar daripada nombor telefon 11 digit biasa hampir pasti penipuan.
  \item Sandaran faktor dua: Jangan tukar sesuka hati nombor telefon yang terikat kepada akaun penting; jika terpaksa tukar nombor, pastikan dahulu nyah-ikat semua akaun penting, jika tidak pemilik baharu boleh menetapkan semula kata laluan anda melalui kod pengesahan SMS.
\end{itemize}

\end{tcbwarn}

\booksection{WeChat/QQ: Pengurusan Aplikasi Sosial Peringkat Nasional}

WeChat dan QQ telah melampaui peranan alat komunikasi, menjadi “Super App” yang menggabungkan sosial, pembayaran dan capaian maklumat.

\booksubsection{Sandaran dan pemindahan rekod perbualan:}

\begin{enumerate}
  \item Log masuk WeChat versi PC atau versi Mac pada komputer.
  \item Klik menu di penjuru kiri bawah > [Pindah dan Sandar] > [Sandar dan Pulih].
  \item Pilih "Sandarkan rekod perbualan ke komputer", kemudian pilih perbualan yang perlu disandarkan pada telefon.
  \item Pindah ke telefon baharu: Pada WeChat telefon lama, [Saya] > [Tetapan] > [Umum] > [Pindahan dan Sandaran Rekod Perbualan] > [Pindah] > [Pindah ke Telefon/Tablet], imbas kod QR untuk selesai.
\end{enumerate}

\booksubsection{Kaedah mendalam pembersihan ruang storan:}

\begin{itemize}
  \item Pintu masuk biasa: WeChat [Saya] > [Tetapan] > [Umum] > [Ruang Storan], boleh membersihkan cache dan membersihkan rekod perbualan secara halus mengikut masa dan jenis fail.
  \item Pembersihan dalam perbualan: Masuk sesebuah perbualan > "…" di penjuru kanan atas > [Rekod Perbualan] > [Urus Ruang Storan], boleh membersihkan gambar, video dan fail dalam perbualan tersebut secara berasingan sambil mengekalkan teks.
  \item Pembersihan gambar/video asal: Gambar dan video HD asal yang diterima dalam kumpulan paling banyak memakan ruang; boleh tandakan "gambar asal" dalam [Urus Ruang Storan] untuk memadam secara pukal.
  \item Matikan muat turun automatik: [Tetapan] > [Umum] > [Foto, Video, Fail dan Panggilan], matikan "Muat Turun Automatik" untuk mengelak cache fail kumpulan yang banyak terkumpul secara senyap di bawah WiFi.
  \item Cache mini program: Dalam [Temui] > [Mini Program], tekan lama mini program yang jarang digunakan untuk memadam, ini boleh melepaskan cache tempatannya.
  \item Sandaran mampat rekod perbualan: Semasa memsandarkan di WeChat versi komputer, tandakan "Mampat"; rekod perbualan berukuran beberapa GB boleh dimampatkan kepada 30\%-50\% daripada saiz asal.
\end{itemize}

\booksubsection{Tip pengurusan kumpulan:}

\begin{itemize}
  \item Pengumuman kumpulan: Pemilik/pentadbir kumpulan menerbitkannya melalui "…" di penjuru kanan atas perbualan kumpulan > [Pengumuman Kumpulan], peringatan kuat kepada semua ahli, sesuai untuk notis penting. Dipaparkan secara automatik apabila ahli baharu menyertai kumpulan.
  \item Tugasan kumpulan: Pemilik kumpulan boleh menetapkan mesej yang disemat sebagai "tugasan kumpulan", semua ahli akan melihat petunjuk merah di bahagian atas skrin, hilang secara automatik selepas klik selesai. Sesuai untuk pendaftaran aktiviti dan kutipan tugasan.
  \item Rantaian kumpulan: Taip "\#Jielong" dalam kotak input untuk memulakan rantaian, sesuai untuk menempah makanan atau mengira pendaftaran. Hasil rantaian dihimpunkan secara automatik di bahagian atas kumpulan, mengelakkan banjir mesej.
  \item Siaran langsung kumpulan: "+" di penjuru kanan atas perbualan kumpulan > [Siaran Langsung Kumpulan], boleh memulakan siaran langsung yang hanya boleh ditonton oleh kumpulan ini, sesuai untuk latihan dalaman syarikat, perjumpaan keluarga dan ujian dalaman produk.
  \item Lipat kumpulan ini: Dalam senarai kumpulan, leret ke kiri nama kumpulan > [Lipat], boleh menyimpan kumpulan yang tidak penting dalam "Kumpulan Dilipat", hanya memberi peringatan apabila ada yang @anda.
  \item Hanya pemilik/pentadbir kumpulan boleh @semua: Apabila ahli kumpulan melebihi 500 orang, ahli biasa tidak boleh @semua orang, untuk mengelakkan gangguan.
\end{itemize}

\booksubsection{Pengurusan halus kebenaran Moments:}

\begin{itemize}
  \item Keterlihatan mengikut kumpulan: Semasa menyiarkan, klik "Siapa boleh lihat", pilih "Sebahagian boleh lihat" atau "Jangan tunjuk kepada", sasarkan dengan tepat berdasarkan kumpulan tag yang ditetapkan lebih awal (contohnya "rakan sekerja", "keluarga", "pelanggan").
  \item Julat masa rakan boleh lihat: [Saya] > [Tetapan] > [Kebenaran Rakan] > [Moments] > [Julat Masa Rakan Boleh Lihat Moments], boleh pilih "3 hari / 1 bulan / 6 bulan / Semua", disyorkan menetapkan "6 bulan" untuk mengimbangi privasi dan kehidupan sosial.
  \item Blok satu siaran Moments: Tekan lama sesuatu siaran dalam Moments > [Jangan lihat Moments dia] atau [Berhenti mengikut], pihak tersebut tidak akan mengetahuinya.
  \item Pengurusan iklan Moments: [Saya] > [Tetapan] > [Mengenai WeChat] > [Maklumat Peribadi dan Kebenaran] > [Pengurusan Iklan Peribadi], boleh mematikan cadangan peribadi berasaskan tingkah laku.
  \item 10 siaran untuk orang asing: [Kebenaran Rakan] > [Moments] > [Benarkan Orang Asing Melihat 10 Siaran Moments], matikannya untuk mengelakkan orang asing mengimbas siaran lampau anda.
\end{itemize}

\booksection{Perbandingan Panggilan Video}

Setiap alat panggilan video ada kekuatannya yang tersendiri; memilih alat yang sesuai mengikut situasi menjadikan komunikasi lebih lancar.

\begin{itemize}
  \item WeChat Video: Liputan peringkat nasional, sesuai untuk semua peringkat umur. Sesuai untuk hubungan harian keluarga dan rakan, komunikasi pantas. Menyokong panggilan video sehingga 20 orang (memerlukan mini program "Video Kumpulan"), kekurangannya kualiti gambar menurun dengan jelas apabila rangkaian lemah.
  \item FaceTime: Eksklusif untuk ekosistem Apple, kualiti gambar cemerlang, penyulitan hujung-ke-hujung. Bermula dengan iOS 15 menyokong jemputan melalui "pautan" untuk pengguna Android/Web menyertai, tetapi hos mesti menggunakan peranti Apple. Sesuai untuk komunikasi keluarga merentas peranti.
  \item Zoom: Pilihan utama mesyuarat antarabangsa, menyokong mesyuarat percuma 40 minit sehingga 100 orang. Fungsi perkongsian skrin, latar maya, bilik perbincangan kumpulan lengkap, sesuai untuk kerjasama merentas negara atau jarak jauh.
  \item Tencent Meeting: Pilihan utama pejabat di China, integrasi mendalam dengan WeChat, masuk dengan imbas kod sahaja. Menyokong mesyuarat 300 orang (versi berbayar), rakaman awan, sari kata masa nyata dan fungsi minit mesyuarat yang berkuasa, sesuai untuk latihan dalaman syarikat dan mesyuarat besar.
  \item Cadangan situasi penggunaan: Untuk hubungan harian 1-lawan-1 pilih WeChat; antara pengguna Apple pilih FaceTime; untuk mesyuarat perniagaan rasmi pilih Zoom atau Tencent Meeting (di China pilih Tencent, merentas sempadan pilih Zoom).
\end{itemize}

\booksubsection{Rujukan keperluan lebar jalur rangkaian dan penggunaan data:}

\begin{itemize}
  \item Panggilan video WeChat: Memerlukan kira-kira 500Kbps-1Mbps lebar jalur, kualiti standard menggunakan kira-kira 3-5MB data seminit, HD kira-kira 8-12MB.
  \item FaceTime: SD kira-kira 1Mbps, 720p kira-kira 2Mbps, 1080p memerlukan 4Mbps ke atas. Data kira-kira 5-15MB seminit.
  \item Zoom/Tencent Meeting: 720p kira-kira 1.5-3Mbps, 1080p memerlukan 3-5Mbps. Data kira-kira 300-800MB sejam, hos mesyuarat berbilang orang menggunakan lebih banyak data.
  \item Cadangan: Untuk panggilan video yang panjang, gunakan WiFi; pada rangkaian mudah alih, turunkan kualiti gambar secara aktif (matikan "video HD" dalam tetapan aplikasi) untuk menjimatkan data.
\end{itemize}

\booksubsection{Tip pengoptimuman kualiti panggilan video:}

\begin{itemize}
  \item Cahaya: Belakangi tingkap, hadapi sumber cahaya (cahaya semula jadi paling baik), elakkan cahaya dari belakang yang menyebabkan wajah menjadi gelap.
  \item Sudut: Kamera telefon sedikit lebih tinggi daripada mata untuk mengelakkan dagu dua; semasa panggilan dalam mod landskap, pegang dengan dua tangan lebih stabil.
  \item Audio: Gunakan fon telinga untuk mengelakkan gema dan bunyi ciapan; dalam persekitaran senyap jangan aktifkan mod "Pembesar Suara".
  \item Rangkaian: Utamakan WiFi; semasa panggilan video 4G/5G cuba berhampiran tingkap untuk mengelakkan penurunan isyarat; dalam mesyuarat berbilang orang boleh matikan video orang lain dan kekalkan hanya video penceramah.
  \item Kaburkan latar: Tencent Meeting dan Zoom kedua-duanya menyokong kaburan AI untuk menutup latar yang bercelaru. Perhatikan memakai pakaian berwarna polos, elakkan warna yang hampir sama dengan latar kerana menyebabkan "kegagalan skrin hijau".
\end{itemize}

\begin{tcbinfo}{Petua Ringkas}
Apabila rangkaian tersekat-sekat, matikan video dan kekalkan suara terlebih dahulu — lebih sopan daripada terus menutup panggilan. Sebelum mesyuarat penting, boleh gunakan fungsi “Semakan Rangkaian” bawaan “Tencent Meeting” untuk menguji lebar jalur lebih awal.

\end{tcbinfo}

\booksection{E-mel: Pemasangan dan Penggunaan pada Telefon}

\begin{itemize}
  \item Tambah akaun e-mel: Pada iOS di Tetapan > Mel > Akaun > Tambah Akaun; Android biasanya mempunyai aplikasi "E-mel".
  \item Tip penggunaan: Tetapkan cara pemberitahuan yang berbeza untuk akaun e-mel yang berbeza; gunakan gerak isyarat leret untuk mengarkib, memadam atau menanda e-mel dengan pantas.
\end{itemize}

\booksubsection{Perbezaan protokol IMAP vs POP3 vs Exchange:}

\begin{itemize}
  \item POP3: Selepas memuat turun mel daripada pelayan, yang asal dipadam secara lalai, hanya disimpan secara tempatan. Sesuai untuk situasi yang hanya menggunakan satu peranti dan memerlukan bacaan luar talian. Kekurangannya, selepas bertukar peranti, mel lama tidak dapat dilihat; sudah tidak disyorkan.
  \item IMAP: Mel sentiasa disimpan di pelayan, secara tempatan hanya sebagai cache. Status segerak konsisten merentas pelbagai peranti (telefon/komputer/tablet) — satu pihak dibaca, semua dibaca. Protokol pilihan arus perdana masa kini.
  \item Exchange (termasuk ActiveSync): Protokol Microsoft; selain mel, ia juga menyegerakkan kalendar, kenalan dan tugasan. Pilihan pertama untuk e-mel korporat. Outlook, NetEase Enterprise Mail dan Tencent Enterprise Mail semuanya menyokongnya.
  \item Cadangan pemilihan: E-mel peribadi gunakan IMAP; e-mel syarikat utamakan Exchange, kemudian IMAP; elakkan POP3 sebanyak mungkin, kecuali ada keperluan arkib yang khas.
\end{itemize}

\booksubsection{Kaedah konfigurasi e-mel korporat:}

\begin{enumerate}
  \item iOS: [Tetapan] > [Mel] > [Akaun] > [Tambah Akaun] > pilih "Microsoft Exchange" atau "Lain-lain".
  \item Android: Dalam aplikasi "E-mel" bawaan atau aplikasi Outlook, pilih "Tambah akaun", masukkan alamat e-mel penuh dan kata laluan.
  \item Pengisian pelayan (contoh Exchange): Alamat pelayan biasanya "mail.namasyarikat.com" atau "outlook.office365.com", bahagian domain diisi dengan singkatan syarikat.
  \item Jika muncul mesej "tidak dapat mengesahkan pelayan", boleh masukkan secara manual: pelayan masuk (IMAP) port 993, penyulitan SSL; pelayan keluar (SMTP) port 465 atau 587, penyulitan TLS.
  \item Selepas log masuk, disyorkan mengaktifkan "Push" dalam [Tetapan Akaun] untuk menerima mel secara masa nyata, dan mematikan "Pratonton Mel" untuk melindungi rahsia perniagaan.
\end{enumerate}

\begin{tcbinfo}{Petua Kecil}
E-mel korporat biasanya memerlukan pentadbir mengaktifkan “kata laluan kebenaran klien” atau “kod kebenaran IMAP/SMTP”, tidak boleh menggunakan terus kata laluan log masuk web. Apabila log masuk gagal, jana dahulu kod kebenaran dalam tetapan mel versi web, kemudian masukkan ke dalam telefon.

\end{tcbinfo}

\begin{tcbpractice}{Amali: Sandarkan sekali rekod perbualan WeChat ke komputer}
\begin{enumerate}
  \item Mulakan sandaran pada komputer: Buka WeChat versi komputer > menu di penjuru kiri bawah > [Pindah dan Sandar] > [Sandar dan Pulih].
  \item Pilih sandar ke komputer: Klik "Sandarkan rekod perbualan ke komputer".
  \item Pilih kandungan sandaran pada telefon, tunggu sehingga sandaran selesai.
\end{enumerate}

\end{tcbpractice}
